\documentclass[12pt]{article}
\usepackage{amsmath, amssymb, geometry, enumitem, booktabs}
\geometry{margin=1in}
\setlength{\parindent}{0pt}

\title{Solutions -- Quiz 4\\ \vspace{8pt} \Large ECON 308 -- International Economic Relations}
\author{Instructor: Muhebullah Karimzada}
\date{October 28,  2025}
\usepackage{comment}

\begin{document}
\maketitle

\section*{Multiple Choice (1 mark each)}

\textbf{1.} \emph{If a country’s terms of trade improve, its welfare will:} \\
\textbf{Correct: A.} \textit{Always increase.} In the standard (distortion-free) model, an improvement in the terms of trade (export price relative to import price) pivots the trade (isovalue) line outward around the production point, allowing the country to reach a higher indifference curve.\\[0.6em]

\textbf{2.} \emph{Growth biased toward a country’s export good tends to:}\\
\textbf{Correct: B.} \textit{Deteriorate its terms of trade.} Export-biased growth expands the world supply of the export good and tends to lower its relative price, worsening the home country’s terms of trade (for a price-taking small country, prices are exogenous, but the sign remains the same in general equilibrium intuition).\\[0.6em]

\textbf{3.} \emph{If the world relative price of the country’s export good rises due to global demand:}\\
\textbf{Correct: A.} \textit{Improve TOT and raise welfare.} A higher relative price of exports is an improvement in the TOT, shifting the consumption possibility outward and increasing welfare.\\[0.6em]

\textbf{4.} \emph{If the world relative price of cloth rises and the country exports cloth, it will:}\\
\textbf{Correct: B.} \textit{Experience an improvement in its terms of trade.} The export good becomes relatively more valuable, raising real income at given production.\\[0.6em]

\textbf{5.} \emph{In the Standard Trade Model, changes in welfare are illustrated by:}\\
\textbf{Correct: C.} \textit{Tangency between the highest indifference curve and the trade line.} Welfare rises when the budget (trade) line at world prices supports a higher indifference curve.

\hrule
\vspace{1em}
\newpage

\subsection*{Problem 1 (7 marks)}
\textbf{Data.} Goods: food ($F$) and manufactures ($M$). PPF: $F^2 + M^2 = 100$. World relative price: $P_M/P_F = 2$. Consumption satisfies $M = 0.5F$.

\paragraph{1) Production point $(F_P,M_P)$ at given prices.}
With $F$ on the horizontal axis and $M$ on the vertical, the PPF slope is the marginal rate of transformation (MRT):
\[
\text{PPF: } F^2 + M^2 = 100 \;\Rightarrow\; 2F\,dF + 2M\,dM = 0 \;\Rightarrow\; \frac{dM}{dF} = -\frac{F}{M}.
\]
The isovalue (trade) line has slope $-\frac{P_F}{P_M}$. Tangency (value-maximizing production) requires
\[
-\frac{F}{M} = -\frac{P_F}{P_M} \quad \Rightarrow \quad \frac{F}{M} = \frac{P_F}{P_M} = \frac{1}{2}\quad\Rightarrow\quad F=\frac{M}{2}.
\]
Impose on PPF:
\[
\left(\frac{M}{2}\right)^2 + M^2 = 100 \;\Rightarrow\; \frac{M^2}{4}+M^2 = \frac{5}{4}M^2=100 \;\Rightarrow\; M^2=80 \;\Rightarrow\; M_P=\sqrt{80}=4\sqrt{5}\approx 8.944.
\]
Hence $F_P=\frac{M_P}{2}\approx 4.472$.
\[
\boxed{(F_P,M_P) \approx (4.472,\;8.944)}.
\]

\paragraph{2) Consumption point $(F_C,M_C)$ given $M=0.5F$ and trade at world prices.}
Normalize prices by $P_F=1,\;P_M=2$ (only relative prices matter). The value of production (national income at world prices) is
\[
V = P_F F_P + P_M M_P = 1\cdot 4.472 + 2\cdot 8.944 \approx 4.472 + 17.888 = 22.36.
\]
Budget (trade) line at world prices:
\[
P_F F + P_M M = 22.36 \quad \Rightarrow \quad F + 2M = 22.36.
\]
Consumption must lie on the ray $M=0.5F$. Substitute into the budget:
\[
F + 2(0.5F) = 2F = 22.36 \;\Rightarrow\; F_C = 11.18,\qquad M_C = 0.5F_C = 5.59.
\]
\[
\boxed{(F_C,M_C) \approx (11.18,\;5.59)}.
\]
\textit{Trade volumes:} Exports of $M$: $M_P - M_C \approx 8.944-5.59 = 3.354$. Imports of $F$: $F_C - F_P \approx 11.18-4.472 = 6.708$. Valued at world prices, export revenue $=2\times 3.354=6.708$ equals import expenditure $=1\times 6.708$ (trade is balanced).

\paragraph{3) Terms of trade and gains from trade.}
The country’s terms of trade are $P_M/P_F = 2$ (given). Production is chosen to maximize the value at these prices; then trading at the same prices allows the country to move from $(F_P,M_P)$ on the PPF to $(F_C,M_C)$ on a \emph{higher} indifference curve (the trade line through the production point supports a strictly higher utility than autarky, unless preferences happen to make the autarky point coincidentally optimal at those prices). Hence, under standard assumptions:
\[
\boxed{\text{Welfare rises with trade at } P_M/P_F=2.}
\]

\bigskip
\hrule
\bigskip

\subsection*{Problem 2 (8 marks)}
\textbf{Data.} Before growth: $Q_T = 200 - 2Q_C$ (textiles $T$, chemicals $C$), world price ratio $P_T/P_C=1$. After biased growth toward $T$: $Q_T = 240 - 2Q_C$.

\paragraph{1) Change in maximum production possibilities (intercepts).}
Intercepts are where the other good is zero.

\textit{Before growth:} If $Q_C=0$, $Q_T^{\max}=200$; if $Q_T=0$, $Q_C^{\max}=100$.\\
\textit{After growth:} If $Q_C=0$, $Q_T^{\max}=240$; if $Q_T=0$, $Q_C^{\max}=120$.

So
\[
\boxed{\Delta Q_T^{\max}=+40,\quad \Delta Q_C^{\max}=+20.}
\]
Growth is $T$-biased (larger outward shift in the $T$-intercept).

\paragraph{2) Change in export potential at fixed world prices ($P_T=P_C=1$).}
At prices $(1,1)$, the isovalue line has slope $-1$ while the PPF slope is constant at $-2$. Value maximization chooses the corner with the flatter price line: set $Q_C=0$ (produce only $T$).
\[
\text{Before: } (Q_T,Q_C)=(200,0)\Rightarrow \text{income }V=200.
\qquad
\text{After: } (Q_T,Q_C)=(240,0)\Rightarrow V=240.
\]
With homothetic preferences (standard in the STM), expenditure shares depend only on prices. Let $s_T$ be the consumption share of $T$ at $(1,1)$. Then exports of $T$ equal production minus consumption:
\[
X_T = Q_T - \frac{s_T V}{P_T} = (1-s_T)Q_T
\]
(since $V=P_T Q_T + P_C Q_C = Q_T$ at $Q_C=0$ and $P_T=P_C=1$). Therefore
\[
X_T^{\text{after}} = (1-s_T)\cdot 240 = \frac{240}{200}\, X_T^{\text{before}}.
\]
\[
\boxed{\text{Export potential of }T\text{ rises by }20\% \text{ at the given prices.}}
\]

\paragraph{3) Effects on terms of trade (TOT) and welfare.}
\begin{itemize}[leftmargin=1.2em]
\item \textbf{Small-country case (prices fixed at world levels):} Income rises from $200$ to $240$ and the consumption set expands; welfare \emph{unambiguously increases}. Terms of trade are unchanged by assumption.
\item \textbf{Large-country case (endogenous prices):} $T$-biased growth increases the world supply of $T$, tending to \emph{lower} $P_T/P_C$ (a TOT \emph{deterioration}). Welfare depends on the balance between higher productive capacity and worse TOT. Typically welfare still rises; only in extreme cases with very elastic foreign demand could TOT deteriorate so much that welfare falls (\textit{immiserizing growth}).
\end{itemize}
\[
\boxed{\text{$T$-biased growth tends to worsen the home TOT.}}
\]

\[
\boxed{\text{Welfare rises unless the TOT deterioration is exceptionally large.}}
\]

\end{document}
